\documentclass[letterpaper]{article} % DO NOT CHANGE THIS
\usepackage[submission]{aaai2027}  % DO NOT CHANGE THIS
\usepackage[hyphens]{url}  % DO NOT CHANGE THIS
\usepackage{graphicx} % DO NOT CHANGE THIS
\urlstyle{rm} % DO NOT CHANGE THIS
\def\UrlFont{\rm}  % DO NOT CHANGE THIS
\usepackage{natbib}  % DO NOT CHANGE THIS
\usepackage{caption} % DO NOT CHANGE THIS
\frenchspacing  % DO NOT CHANGE THIS
\usepackage{booktabs}
\usepackage{amsmath,amssymb}
\usepackage{tikz}
\usetikzlibrary{shapes.geometric,arrows.meta,positioning}
\pdfinfo{
/TemplateVersion (2027.1)
}
\setcounter{secnumdepth}{0}

\title{Opportunity Certification: Deciding Whether a Learned Component\\ Is Warranted, Before Training One}
\author{Anonymous Submission}
\affiliations{}

\newcommand{\x}{$\times$}
\newcommand{\orc}{$^{\dagger}$} % marks a clairvoyant (non-deployable) ingredient
\newcommand{\rung}[1]{\textsc{#1}}
% Use the external figure when it is present, otherwise a visible placeholder, so the
% document still compiles on a checkout that lacks the exported artwork.
\newcommand{\figorbox}[2]{%
  \IfFileExists{#1.pdf}{\includegraphics[#2]{#1.pdf}}{%
  \IfFileExists{#1.png}{\includegraphics[#2]{#1.png}}{%
  \fbox{\parbox[c][2.4cm][c]{0.9\columnwidth}{\centering\ttfamily\small [#1 not found]}}}}}

\begin{document}
\maketitle

\begin{abstract}
\noindent A learned model is warranted when an oracle's advantage is \emph{reachable} without the oracle's information, and a signal that a model can rank correctly may still be one it cannot act on: \textbf{rank-learnable is not placeable}. This paper turns both observations into a procedure, \textbf{Opportunity Certification}, that decides whether to build a learned component \emph{before} one is trained, and demonstrates it across four settings. Stage~1, the \emph{Instrument}, asks how much of a measured oracle gap is real, comparing an adversarially tuned baseline against an oracle held to the same bandwidth and restricted to information a deployable predictor could hold. Stage~2, the \emph{Necessity Ladder}, asks whether a real gap is reachable, granting an oracle one perfect capability at a time; if perfection in a capability cannot close the gap, no learner improving that capability can. Instantiated in cache prefetching, Stage~1 finds 14--87\% of the apparent opportunity phantom across twelve production traces, surviving on three at 12.7--19.7 points; Stage~2 then finds every deployable capability insufficient, for one measured reason, \emph{endogenous slack}: a policy's own mistimed fetches evict its future targets and destroy the residence window its timing model must hit, so a hundred-request survival probability falls from 1.0 to 0.05. Applied unchanged, the procedure returns no corridor for language-model cache warming and \emph{build} for learned indexes on two of five key sets, precisely where the oracle's advantage rests on the data rather than the future. Four settings, four pre-registered verdicts, none requiring a trained model to reach. The procedure, not any single number, is the contribution.
\end{abstract}

\section{Introduction}
A recurring reflex in machine learning for systems is to measure the gap between a deployed heuristic and an oracle and to treat that gap as an opportunity a model should be trained to capture. The inference has two silent premises: that the gap was measured honestly, and that the information the oracle used is information a deployable model could hold. This paper argues that both premises fail often enough that the inference itself should be tested, and tested \emph{before} a model is trained. The argument is developed in cache prefetching, where both failures occur and can be measured, and the resulting procedure is then applied unchanged to two further domains.

The prevailing yardsticks invite mismeasurement. The gap to Belady's optimum \citep{belady} ignores prefetching, and its prefetch-aware repair \citep{demandmin} shows that the optimum itself moves once prefetching is admitted. Oracle ceilings spend unbounded bandwidth against throttled or under-tuned baselines, and clairvoyant ceilings credit \emph{compulsory} misses, first references to never-seen objects that no history-based policy can prefetch. In preliminary measurement for this work, a $+42$-point synthetic headroom collapsed to $-0.02$ once the baseline was tuned across a predictor family and the oracle was charged the baseline's exact prefetch bandwidth. That collapse fixed the two commitments everything below enforces: a strongest-available baseline, and every ceiling held to the baseline's byte budget, under pre-registered thresholds and machine-checked invariants.

The proposal, \textbf{Opportunity Certification}, submits an apparent oracle gap to two stages (Figure~\ref{fig:overview}). \emph{Stage 1, the Instrument,} certifies how much of the gap is real, measuring a \textbf{learnable corridor}: the distance between the best decoupled predictor over a swept family and a ceiling restricted to objects a history-based predictor could name, both at one prefetch byte budget. \emph{Stage 2, the Necessity Ladder,} certifies whether a real corridor is reachable, pricing each capability a learner could contribute with a separate oracle. The output is a certificate: \emph{no corridor}, a \emph{trap}, or \emph{build}.

On twelve production traces the Instrument deflates gross corridors by 14--87\%; three survive, at $+12.7$ to $+19.7$ hit-rate points with bootstrap confidence intervals clear of a pre-registered bar. The Ladder then prices arbitration, coverage, and timing on the survivors, and every deployable ingredient falls short. The failure has a single measured mechanism, \textbf{endogenous slack}: a policy's own mistimed fetch volume evicts its future targets, shrinking the residence window its timing model must hit, so a reuse signal can be \emph{rank-learnable} without being \emph{placeable}. A policy built specifically against this bottleneck also fails, because its added signal proves redundant with the predictor's own confidence.

Two further instantiations show the procedure discriminates rather than merely deflates. On key-value cache warming in LLM serving it returns \emph{no corridor}, because a tuned association table already sits on the physically reachable ceiling. On learned indexes it returns \emph{build} on two of five key sets, because there the oracle's advantage rests on the key distribution rather than the future, and it withholds the verdict on the remaining three, where a tuned classical structure already attains that ceiling. Reachability separates a build from a trap; baseline saturation separates a build from an illusion.

\paragraph{Contributions.} (i) Opportunity Certification, a two-stage, domain-general procedure with pre-registered thresholds and machine-checkable invariants; (ii) the Instrument and a deflationary measurement across 12 traces, with a taxonomy of four inflation mechanisms and the finding that eviction and prefetching are substitutes; (iii) the Necessity Ladder and its threefold negative on the surviving corridors; (iv) the endogenous-slack mechanism, separating rank-learnability from placement; (v) a cross-domain instantiation returning four distinct verdicts, with a captured corridor as a positive control and a trained-model competence control establishing that the substrate admits learned wins.

\begin{figure}[t]\centering
\figorbox{Figures/overview}{width=\columnwidth}
\caption{\textbf{Opportunity Certification.} Stage 1 (the Instrument) certifies how much of an apparent oracle gap is real; Stage 2 (the Necessity Ladder) certifies whether a deployable policy can reach it. The certificate reads no corridor, trap, or build.}
\label{fig:overview}
\end{figure}

\section{Related Work}
\textbf{Optimal replacement and prefetch-aware repairs.} Belady's MIN \citep{belady} is optimal for demand caching without prefetching; \citet{demandmin} showed that under prefetching the optimum becomes bandwidth-relative. This work transports that insight to variable-size object caches and adds the reachability split: it bounds what a history-based policy could do, not what an omniscient one could. \citet{cao} treated integrated prefetching and caching under clairvoyance, uniform sizes, and a latency model, assumptions that fail for CDN and key-value objects, so ceilings here are measured per trace.

\textbf{Learned caching.} LRB \citep{lrb} imitates a relaxed Belady oracle for eviction; Cold-RL \citep{coldrl}, DEAP \citep{deap}, DeePref \citep{deepref}, and LightCacheRL \citep{lightcache} learn eviction or prefetching; Pythia \citep{pythia} tunes hardware prefetching with online RL; \citet{jointhw} share representations across replacement and prefetching. \textbf{Baleen} \citep{baleen} is the nearest neighbor and its reported gains are not disputed, because the measured quantity differs on four axes: it optimises disk-head time under a flash write budget rather than hit ratio at a tuned baseline's byte rate; its dominant lever is admission, so the endogenous-slack loop is largely absent; its OPT-derived bound is not restricted to a predictor's vocabulary (M4); and it is compared against production configurations rather than a family-swept baseline (M1). A gain there and a trap here can both be true. To current knowledge no prior work constructs a pre-registered, bandwidth-matched, learnability-split headroom measurement for object caches, nor gates a learned component on an oracle decomposition of necessity.

\textbf{LLM serving caches and learned indexes.} Mooncake \citep{mooncake} released the production KV-cache traces used below; PagedAttention \citep{vllm} and CacheGen \citep{cachegen} established and compressed prefix-shared KV caching. That literature reports achieved hit rates; the cross-domain section measures headroom. Learned indexes \citep{kraska,sosd} supply the positive foil. \citet{sober} showed claimed reasoning gains shrink under variance-aware protocols; this paper is the systems analog, adding a capturability test ahead of any training. S3-FIFO \citep{s3fifo} is the fixed evictor throughout, so every corridor is attributable to prefetch scheduling.

\section{Stage 1: The Instrument}
The Instrument is one deterministic replay in which four arms share a single prefetch byte budget, and every degree of freedom that could inflate the corridor is fixed before any real-trace run.

\textbf{Setup.} Each run replays a two-million-request prefix of a production trace through the same cache: capacity 1\% of the trace footprint, a 5\% warmup excluded from scoring, request-weighted hit ratio as the score, S3-FIFO as the evictor in every arm. Predictors are decoupled from the future: they train on the first half of the trace and are frozen. A prefetched object later requested counts as a hit, one evicted before use counts as waste, and every prefetched byte counts as traffic; the accounting is unit-tested.

\textbf{The four arms} (Figure~\ref{fig:instrument}). \emph{BASE} is S3-FIFO with no prefetch; its origin traffic is the denominator. \emph{TUNED BASELINE} is the best non-learned system constructible against the corridor's own interest: the highest hit ratio over $\{$Markov-1, Markov-2, Markov-3, LSTM$\}$ across all confidence thresholds and fan-outs, capped at $1.15\times$ BASE traffic. \emph{GROSS CEILING} is a clairvoyant prefetcher metered by a token bucket to the baseline's byte rate, reported only as a diagnostic because it fetches objects no real predictor could know. \emph{WARM CEILING}, the primary arm, is the same clairvoyant restricted to objects seen during training, at the same budget; whatever the gross oracle wastes on unreachable objects, the warm ceiling re-spends on reachable ones. It is the honest upper bound for a history-based scheduler: perfect timing, realistic knowledge, equal cost.

\textbf{Why this interval.} The corridor runs from the tuned baseline to the warm ceiling because headroom is meaningful only above the strongest system deployable without new ideas, and only below what a history-based policy could in principle reach at equal budget. The tempting shortcut, learnable $\approx$ gross corridor minus cold slice, is wrong in both directions: on Wikipedia 27\% of the gross oracle's useful prefetches were cold, and redirecting that budget buys $+5.4$ warm points, so subtraction misstates the corridor by up to 5.5 points, enough to flip a verdict. Reachability must be priced inside the replay. An invariant enforces it, requiring exactly zero cold prefetch hits from the baseline and warm arms; an earlier ``cold'' definition that violated this invariant, crediting the baseline with 442{,}188 impossible cold hits, was caught by it and replaced.

\textbf{Pre-registration.} The verdict bar, learnable corridor $\ge 8$ points \emph{and} Pareto dominance at the warm ceiling, was fixed before any real-trace run. Baselines escalate in pre-registered stages (coarse Markov grid, then fine thresholds plus Markov-3, then the LSTM), and a live verdict is provisional until it survives the strongest stage; escalation can only raise the bar, so survival is meaningful and death is final.

\begin{figure*}[t]\centering
\figorbox{Figures/instrument}{width=0.95\columnwidth}
\caption{The Instrument: four arms at one shared prefetch byte budget. The tuned baseline and the warm ceiling (clairvoyant, restricted to training-vocabulary objects, re-spending the cold budget) bracket the learnable corridor; the gross ceiling is a diagnostic that over-counts compulsory misses. A trace is live only if the corridor clears 8 points and the warm ceiling Pareto-dominates.}
\label{fig:instrument}
\end{figure*}

\begin{table*}[t]
\centering\small
\setlength{\tabcolsep}{4.5pt}
\begin{tabular}{llccccp{5.2cm}}
\toprule
trace & family & base & tuned baseline & gross corr. & verdict & reason for verdict \\
\midrule
\texttt{wiki\_2019t} & CDN & 0.152 & 0.5531 @1.09\x & $+26.49$ & retained & (none)  \\
\texttt{meta\_rprn} & CDN & 0.390 & 0.6738 @1.03\x & $+32.57$ & removed &  cold-dominated \\
\texttt{meta\_reag} & CDN & 0.538 & 0.7547 @1.11\x & $+24.51^{\ddagger}$ & removed &  cold-dominated \\
\texttt{cluster50} & KV & 0.416 & 0.7080 @1.02\x & $+20.37$ & retained & (none)  \\
\texttt{cluster53} & KV & 0.595 & 0.6101 @1.10\x & $+25.72^{\dagger}$ & retained & (none)  \\
\texttt{cluster26} & KV & 0.786 & 0.8907 @1.09\x & $+10.27$ & \textbf{removed} & traffic-bought (ceiling $1.16\times$ vs bar $1.09\times$) \\
\texttt{cluster10} & KV & 0.500 & 0.7368 @1.00\x & $-1.18$ & removed & degenerate predictability (baseline $>$ ceiling) \\
\texttt{msr\_proj\_0} & block & 0.552 & 0.7568 @1.13\x & $+4.16$ & removed & baseline saturation (robust to any cap) \\
\texttt{msr\_hm\_0} & block & 0.657 & 0.8380 @1.13\x & $+10.08$ & removed &  cold $+$ saturation \\
\texttt{msr\_web\_2} & block & 0.010 & 0.8284 @1.00\x & $-3.81$ & removed & baseline saturation (prec 0.996) \\
\texttt{w105} & block & 0.269 & 0.8494 @1.05\x & $-0.45$ & removed & baseline saturation \\
\texttt{w87} & block & 0.095 & 0.5383 @1.12\x & $+32.67$ & removed$^{*}$ & Pareto fail by $0.01\times$ (boundary) \\
\bottomrule
\end{tabular}
\caption{Gate A across 12 traces. Escalating the baseline moved \texttt{cluster26} below the threshold, and the LSTM recovered no corridor; survival is not a property of the workload family. $^{*}$\texttt{w87} fails Pareto dominance by $0.01\times$, inside display resolution; the pre-registered rule removes it and it is reported rather than dropped, and no claim depends on it. $^{\dagger}$\texttt{cluster53}'s gross corridor is quoted here against its \emph{strongest} bar, the LSTM (0.6101@$1.10\times$); Table~\ref{tab:learnable} decomposes it against its Markov bar (0.5973@$1.01\times$), giving $+24.46$. Both are reported because the LSTM bar is the one that must be beaten, while the decomposition requires the arm whose in-vocabulary set is defined; the verdict is identical either way. $^{\ddagger}$\texttt{meta\_reag} differs by $0.04$ between tables (a re-run at identical settings); the smaller value is used wherever a corridor is claimed.}
\label{tab:gatea}
\end{table*}

\section{Results: Measurement}
\textbf{Which gross corridors survive.} Table~\ref{tab:gatea} reports Gate A. Three points stand out. The tuned baseline is decisive: \texttt{cluster26} falls below the threshold once the baseline escalates. The LSTM never recovers a corridor, so the survivors are not artifacts of a weak predictor. And survival is not a family property: block traces yield both $+10.08$ and $-3.81$, key-value traces both $+25.72$ and $-1.18$.

\textbf{14--87\% of the gross is phantom.} Table~\ref{tab:learnable} decomposes the six gross-live traces. The three with the most impressive gross corridors lose 51--87\% to the cold slice, because high object churn spends the clairvoyant budget on compulsory misses; the two low-churn KV traces lose only 14--20\%. Wikipedia is the instructive middle case: 52\% phantom, yet enough exploitable warm reuse that the honest corridor still clears the bar, but only once the warm ceiling re-spends the cold budget; naive subtraction would report $+7.27$ and falsely kill it. Corridors are paired mean differences over the same post-warmup requests, resampled as 1{,}000 contiguous blocks with 10{,}000 resamples (block, not i.i.d., because hits are autocorrelated through cache state). All three surviving lower bounds clear 8. (\texttt{cluster53}'s decomposition uses its Markov baseline; against its higher LSTM baseline the corridor is $\ge +18.4$, so the verdict stands.)

\begin{table}[t]
\centering\small
\setlength{\tabcolsep}{3.5pt}
\begin{tabular}{lcccc}
\toprule
trace & gross & cold & \% phantom & \textbf{learnable} (95\% CI) \\
\midrule
\texttt{cluster53} & $+24.46$ & 4.62 & 20\% & $\mathbf{+19.67}$ $[19.4,19.9]$ \\
\texttt{cluster50} & $+20.37$ & 2.69 & 14\% & $\mathbf{+17.56}$ $[15.8,19.4]$ \\
\texttt{wiki\_2019t} & $+26.49$ & 19.22 & 52\% & $\mathbf{+12.72}$ $[10.9,14.5]$ \\
\texttt{msr\_hm\_0} & $+10.08$ & 5.82 & 51\% & $+4.89$ (removed) \\
\texttt{meta\_rprn} & $+32.57$ & 27.90 & 87\% & $+4.32$ (removed) \\
\texttt{meta\_reag} & $+24.55$ & 20.71 & 85\% & $+3.63$ (removed) \\
\bottomrule
\end{tabular}
\caption{Warm/cold decomposition on the gross-live traces (every run verifies bar-cold $=0$ and warm-cold $=0$). Pre-registered rule: reliably live only if the CI lower bound $\ge 8$.}
\label{tab:learnable}
\end{table}


\textbf{The verdicts do not depend on the exact bar.} Pre-registration is a process claim a reader cannot audit, so it is supplemented with a robustness result that can be: sweeping the threshold from 4 to 12 points, applied as the rule applies it, to the bootstrap lower bound.

\vspace{2pt}\noindent{\small
\begin{tabular}{lcccccc}
\toprule
trace & CI lo & 4 & 6 & \textbf{8} & 10 & 12 \\
\midrule
\texttt{cluster53} & 19.4 & live & live & \textbf{live} & live & live \\
\texttt{cluster50} & 15.8 & live & live & \textbf{live} & live & live \\
\texttt{wiki\_2019t} & 10.9 & live & live & \textbf{live} & live & dead \\
\texttt{msr\_hm\_0} & 4.89$^{*}$ & live & dead & \textbf{dead} & dead & dead \\
\texttt{meta\_rprn} & 4.32$^{*}$ & live & dead & \textbf{dead} & dead & dead \\
\texttt{meta\_reag} & 3.63$^{*}$ & dead & dead & \textbf{dead} & dead & dead \\
\bottomrule
\end{tabular}}\vspace{2pt}

\noindent($^{*}$point estimates, conservative here since a bootstrap bound could only be smaller.) Every survivor clears 10 and every removal fails at 6. The other thresholds have wider margins still: \rung{r2}'s arbitration gap measured $\pm0.01$ against a 2-point requirement and so survives any bar above $0.1$, while in the index domain \texttt{meta\_rprn} would flip at 51\%, which is why it is reported as marginal. Whether the thresholds were fixed in advance therefore decides no verdict except that one.

\textbf{Four inflation mechanisms, caught on the record.} The deflation decomposes into four mechanisms (Table~\ref{tab:inflation}), each caught by a named component and each first exposed on this work's own results rather than on the literature's. Stage~2 adds five further self-caught errors, two inside the oracle arms themselves, and the sharpest instance, which both removed and restored verdicts under one rule, is reported with the learned-index results below. A protocol's credibility rests on what it catches, including from its own authors; a test that never overturns anything provides no evidence.

\begin{table}[t]
\centering\small
\setlength{\tabcolsep}{3pt}
\begin{tabular}{p{1.55cm}p{2.5cm}p{2.7cm}}
\toprule
mechanism & exhibited by & caught by \\
\midrule
under-tuned baseline & synthetic $+42\to-0.02$; \texttt{cluster26} $+12.5\to$ below bar & family sweep over $\{$Markov-1/2/3, LSTM$\}\times\tau\times k$ \\
bandwidth mismatch & synthetic oracle at $1.27\times$ baseline traffic & token bucket at the baseline's byte rate \\
traffic-bought corridor & \texttt{cluster26}: ceiling $1.16\times$ vs bar $1.09\times$ & Pareto-dominance check \\
cold conflation & \texttt{meta\_rprn} $+32.6$ gross $\to +4.3$ learnable & warm (in-vocabulary) ceiling restriction \\
\bottomrule
\end{tabular}
\caption{Four ways an apparent corridor inflates, each with the arm that removes it.}
\label{tab:inflation}
\end{table}

\textbf{What separates live from dead.} Where no corridor survives, the tuned baseline already reaches the warm ceiling (\texttt{meta\_reag} 0.754 versus 0.791), so the residue is compulsory-cold and unreachable by construction; where one survives, in-vocabulary reuse exists that the tuned predictor fails to schedule (Wikipedia 0.553 versus 0.680), which is exactly the quantity a learned scheduler would target.

\textbf{Eviction and prefetching are substitutes.} A $2\times2$ of $\{$S3-FIFO, Belady$\}\times\{$no prefetch, clairvoyant at matched bandwidth$\}$ measures their interaction. Six traces, all negative: $-1.34$ to $-7.94$. Better prefetching reduces what better eviction can add on every trace measured, so the decoupled design is the indicated one, which justifies freezing S3-FIFO throughout Stage~2.

\textbf{Robustness.} The corridor shrinks monotonically as the cache grows, from $+15.2$ at a 0.1\% cache to $+12.7$ at 1\% and $+5.1$ at 10\% on Wikipedia, so survival is a property of the operating point: two further traces clear the bar at 0.1\% and none at 10\%. Byte-weighting deflates it again, often below zero where the request-weighted figure is large (\texttt{meta\_reag} reads $+26.7$ by request and $+4.9$ by byte at 0.1\%). The surviving corridor therefore belongs to the request-weighted objective that latency-serving tiers are operated against, not to egress cost, and an operator optimizing bytes should read these traces as having no corridor at all. Stage~2's verdicts are unaffected, being negative under the more favourable weighting.


\section{Stage 2: The Necessity Ladder}
Three certified corridors invite the obvious next act, a learned scheduler. The Ladder replaces that reflex with a question sequence: each rung is a pre-registered, replay-only experiment that isolates one capability with an oracle and prices it before any model is trained (Figure~\ref{fig:ladder}, Table~\ref{tab:decomp}). Every arm holds S3-FIFO fixed, respects the baseline's byte budget through the same token bucket, and reports paired block-bootstrap intervals; oracle arms bound what is achievable and are never reported as deployable results. Each rung grants strictly less foresight than the one before, so the sequence walks from reachable in principle to deployable in fact.

\textbf{\rung{r1} (timing ceiling).} An oracle controlling only fetch timing, restricted to objects the predictor could name, recovers 95.8\% of the corridor on Wikipedia and 91.2\% on \texttt{cluster50}: these corridors are questions of \emph{when}, not \emph{which}. On \texttt{cluster53} the same oracle recovers 7.3\%, so its corridor is an object-choice problem; that trace exits the capture study on this pre-registered criterion. \emph{The capture study below therefore runs on two traces, not three}, which is a real limit on the generality of its negative and is stated here rather than left to the scope section. What supports the conclusion beyond trace count is that each rung fails for a separately measured reason and that the failures are mechanistically linked.

\textbf{\rung{r2} (arbitration).} Learning would help arbitration only if choosing fetches as a set beat choosing them independently. Given identical value estimates and budget, an optimal joint chooser beats a per-candidate threshold rule by $+0.01$ and $-0.01$ points, far below the pre-registered 2-point requirement. The choice is separable; a threshold rule suffices, and reinforcement learning over this decision is unnecessary.

\textbf{\rung{r3} (the offered stream).} Perfect timing applied only to the candidates the baseline actually emits lands \emph{below} the baseline: $-7.01$ $[-7.68,-6.33]$ and $-10.82$ $[-11.70,-9.93]$. No scheduler wins by re-timing what is already offered, so the corridor is a problem of \emph{coverage}.

\textbf{\rung{r4} (coverage).} Widening the emission (fan-out 32, no confidence floor) and inserting just before use only those reuses the wide predictor could name in time clears the threshold: $+10.41$ $[8.74,12.08]$ and $+11.42$ $[9.95,12.91]$. The corridor is reachable in principle; breadth supplies candidates and perfectly timed insertion converts them.

\textbf{\rung{r5} (emission-time fetch).} The natural causal version of \rung{r4} fetches each candidate when predicted and holds it until use. Granting this design a \emph{perfect} forecaster and \rung{r4}'s own selection rule, so that the only difference is when the fetch happens, it still loses 10 to 15 points to \rung{r4} in every configuration, capturing at most 6\% of the corridor. Since a perfect forecaster fails when fetching at prediction time, no better forecaster can rescue the design; the missing ingredient is just-in-time insertion.

\textbf{\rung{r6} (causal placement).} \rung{r6} keeps \rung{r4}'s clairvoyant coverage and selection and swaps only the clock: the \texttt{oracle} arm fetches at the true next use and reproduces \rung{r4} ($+9.69$, $+11.69$, a passing construction check); the \texttt{model} arm fetches at a time predicted causally from recency, frequency, and confidence. The model passes the standard learnability gate (rank correlation $0.60$/$0.43$, never-reused AUC $0.92$/$0.81$) yet lands $-40.06$ $[-42.7,-37.4]$ and $-26.01$ $[-27.9,-24.1]$: precision collapses to $0.05$ and volume grows to $1.5$M fetches against the baseline's $1.0$M. A pre-registered retrain on the wide stream raises rank correlation to $0.68$ without improving capture.

\begin{figure}[t]\centering
\figorbox{Figures/ladder}{width=\columnwidth, height=7.6cm}
\caption{The Necessity Ladder as run (wiki / cluster50). Each rung is a pre-registered oracle decomposition pricing one ingredient; green rungs certify the corridor reachable in principle (\rung{r1} timing, \rung{r4} coverage), red rungs are negatives. \rung{r6}'s failure is explained by the mechanism band beneath it, and \rung{r7} is a policy built specifically to escape that mechanism. The certificate reads \emph{trap}: real, but reachable only with foresight.}
\label{fig:ladder}
\end{figure}

\begin{table}[t]
\centering\small
\setlength{\tabcolsep}{3pt}
\begin{tabular}{lllcc}
\toprule
selection & insertion & rung & wiki & c50 \\
\midrule
wide\orc & clairvoyant JIT\orc & \rung{r4} coverage & $+10.41$ & $+11.42$ \\
$k{=}1$ stream & clairvoyant JIT\orc & \rung{r3} stream & $-7.01$ & $-10.82$ \\
wide\orc & at emission & \rung{r5} emission & $+0.62$ & $-3.07$ \\
wide\orc & causal hazard & \rung{r6} placement & $-40.06$ & $-26.01$ \\
armed wide & causal event & \rung{r7} CGP & $-7.60$ & $-5.39$ \\
causal & causal & deployable & \multicolumn{2}{c}{$\le\min$ above} \\
\bottomrule
\end{tabular}
\caption{Capture decomposition (corridor versus baseline, points). \orc\ marks a clairvoyant ingredient. Every mixed row that removes one collapses; the all-causal rows sit far below the baseline. \rung{r5} is quoted in its \emph{most favourable} configuration, a perfect forecaster under a cumulative rather than instantaneous budget, so it is an upper bound on that design. \rung{r1} and \rung{r2} are reported in text.}
\label{tab:decomp}
\end{table}

\textbf{The mechanism: endogenous slack.} Why does a model that ranks reuse well still fail to place fetches? Under the baseline's light load a prefetched object nearly always survives until use; the hundred-request survival probability is essentially $1$, an apparently wide window. That window, however, is set by the policy's own behaviour, not the workload. A wide causal policy that mistimes issues far more fetches, its precision falls toward $0.05$, and the extra admissions evict useful objects sooner, shrinking the very window each later fetch must hit: measured directly, the probability that a prefetched object is still cached one hundred requests after admission falls from $1.00$ under the baseline to $0.055$ on Wikipedia and $0.135$ on \texttt{cluster50}. Concretely, a correctly ranked and reasonably timed fetch of an object $C$ converts under baseline load because little else displaces it; the same fetch under the policy's own flood of speculative admissions finds $C$ evicted before its request arrives. The policy did not misjudge $C$; its own volume destroyed the placement. Hence the general lesson: \textbf{a signal can be rank-learnable without being placeable}. A learnability test certifies only ordering, whereas placement requires landing inside a window the decision itself narrows; to current knowledge this is the first direct demonstration that the gap between the two can be complete.

\textbf{\rung{r7}: the constructive escape also fails.} Endogenous slack suggests its own remedy, coverage without volume. Corroboration-Gated Prefetch (CGP) arms wide-predictor candidates onto a byte-free watchlist and fires an actual fetch only when a tight predictor's signal triggers, choosing among watchlisted candidates by confidence within the baseline's exact budget. Coverage comes from the free watchlist, volume stays at the baseline's, and with perfect components the design reduces to \rung{r4}. It fails regardless: $-7.60$ $[-8.1,-7.1]$ and $-5.39$ $[-5.9,-4.9]$. The reason is specific: raising the watchlist weight shifts the funded set from 43\% to 69\% watchlisted while precision stays at $0.268$, so the gate adds $+0.18$ and $-0.04$ points. The corroboration signal is statistically redundant with the predictor's own confidence. No causal policy in this design space widens coverage without widening the low-precision volume that endogenous slack penalises, a stronger negative than a simple failure because the policy was built against the exact bottleneck and the bottleneck absorbed it.

\textbf{What the ladder does and does not establish.} The bottom row of Table~\ref{tab:decomp}, that a deployable policy is bounded by the minimum above it, holds \emph{if} \rung{r3}--\rung{r7} bracket the deployable design space of selection and insertion. That is an inference from named instances plus a mechanism, not an impossibility proof, and it is stated as such: a design outside the decomposition, for instance one that modifies the evictor to protect warmed objects, is not excluded by these measurements, though the substitutes result argues against it. What the rungs do establish is stronger than a single trained model's failure would be, because each removes one clairvoyant ingredient at a time and the collapse is attributable to the ingredient removed.

\textbf{What would have produced a build.} The Ladder was prepared to certify the opposite outcome: a build required one rung to hold under causal information. \rung{r4} shows the target exists and \rung{r6}'s oracle arm shows the scheduler design is sound, so a causal placement arm landing anywhere above the baseline would have read build, and a 2-point gap at \rung{r2} would have recommended a learned arbiter specifically. Neither occurred. The learned-index instantiation below exercises exactly this path and does return build, which is what distinguishes a procedure from a prior.

\section{The Procedure Across Domains}
A measurement procedure earns trust by returning different answers on different inputs. The same arms and bars, unchanged in logic, were carried to two further settings; one returns a third distinct negative, the other a positive, and trained models then check both directions.

\textbf{LLM KV-cache serving: no corridor.} On Mooncake production traces \citep{mooncake} (one hour each of conversation and tool/agent traffic; objects are 512-token KV blocks), warming a block before next use is the prefetch analog and hit rate measures avoided prefill, the time-to-first-token lever. Two adaptations keep the port honest: warming is decided only at request boundaries, since per-access warming would fake-convert whole prefix chains at zero latency, and reachability is physical, since a block never computed cannot be warmed by any system. The verdict holds under both comparators: against the tuned baseline the corridor is negative outright ($-0.56$, $-1.87$), and against the capped-baseline control, the comparison more favourable to a corridor, the margin is $+2.98$ and $+2.14$ versus an 8-point gate. The reason is structural: prefix-chain reuse is near-deterministic (held-out coverage 0.728 versus 0.008 for popularity), so a tuned association table already sits on the physically reachable ceiling. Reuse timing is learnable on the conversation trace ($\rho{=}0.27$, AUC 0.81), but no corridor exists for that signal to capture.

\textbf{Learned indexes: a captured corridor.} In a static learned index the oracle is a perfect index at one probe, and the deployable arm is a two-stage recursive model (RMI) trained on the keys alone, finishing with a last-mile search over its guaranteed error window \citep{kraska}. The decisive difference from caching is reachability: a key's position is $N$ times its cumulative distribution, learnable from the keys themselves, so the oracle's advantage rests on the data rather than the future and the procedure should be able to certify it.

Certification requires sweeping both sides. Binary search is not the tuned non-learned baseline: a cache-conscious B-tree \citep{ccbtree} costs $\lceil\log_{B} N\rceil$ node accesses, and interpolation search \citep{interp} exploits precisely the position-is-CDF signal with no model at all, in $O(\log\log N)$ probes on smooth distributions. The baseline is therefore the best of $\{$binary, B-tree, interpolation$\}$ per dataset; symmetrically, the RMI's leaf count is swept, capped at $N/16$ so the model stays space-comparable to the B-tree rather than degenerating into a lookup table. Both cost units are reported, memory probes and distinct cache lines touched, because they rank these structures differently; the cache-line model treats the RMI root as resident and charges the leaf model and last-mile search.

Table~\ref{tab:index} gives the verdicts: \textbf{build on two of five real key sets}, with sharp discrimination among the rest. A tuned RMI beats the tuned family on \texttt{books} (67.5\% of the corridor by probes, 64.3\% by lines) and on the \texttt{meta\_rprn} identifiers (50.4\%/54.9\%, a thin margin reported as such); it falls short on \texttt{osm\_cellids} and fails outright on \texttt{fb} and the Wikipedia identifiers. One measured quantity explains every verdict: how each leaf's \emph{maximum} error, which sets the search window, responds to capacity. Where the bulk fits within a cache line the corridor is captured (\texttt{meta\_rprn}: 97.5\% of keys inside one line; \texttt{books}: bimodal, 93.6\% inside one line, 6.4\% beyond sixty-four, nothing between). Where capacity helps but runs out, the corridor is approached but missed (\texttt{osm\_cellids}: 13.13 to 7.05 probes as leaves grow, tail share 94\% to 16\%, never reaching the B-tree's 6.00). Where the model is uniformly unfit, capacity is inert: \texttt{fb} and \texttt{wiki\_2019t} hold 100\% of keys beyond sixty-four lines and do not move across a 125-fold leaf increase, a structural failure of piecewise-linear approximation that no capacity budget addresses. The two units also disagree in both directions across the suite (on \texttt{books} probes flatter the model, on synthetic uniform keys cache lines do), so a single-metric evaluation would have returned a wrong verdict in both cases; a correctness invariant, every key inside its declared window, holds on every run.

\textbf{The sweep rule reverses verdicts in both directions.} This suite is where the discipline is most visible, because applying one rule symmetrically both destroyed and restored conclusions. Scoring the learned index against binary search alone, the field's customary comparison, is mechanism M1 exactly; adding the two competitors the rule requires removed up to 337 points of apparent capture and reversed three of five verdicts. Applying the same rule to the learned side then reversed two of them back, because sweeping the baseline while pinning the RMI at one leaf count is the identical error inverted: at fixed capacity \texttt{books} reads 41.7\% and \texttt{meta\_rprn} 10.4\%, both below bar, and once the model is swept as well they read 67.5\% and 50.4\%, both above it. A procedure that only ever removed positives would be a rejection machine; one that removes and restores under a single rule is a measurement.

\textbf{On the number of positives.} Across four settings and eight key sets the procedure returns one unambiguous build and one marginal, and a reader may reasonably ask whether a bar returning so few positives is calibrated. The property under test is not the rate of positives but whether the verdict tracks a measured property of the input, and it does: capture is predicted by how each leaf's maximum error responds to capacity, a quantity measured independently of the verdict and reported in full above, which separates the captured, the approached, and the structurally unfit without reference to the outcome. A bar tuned to produce positives would not exhibit that correspondence.

\begin{table}[t]
\centering\small
\setlength{\tabcolsep}{4pt}
\begin{tabular}{llrrrl}
\toprule
key source & keys & base & probes & lines & verdict \\
\midrule
SOSD \texttt{books} & 2M & btree & 67.5 & 64.3 & \textbf{build} \\
IDs: \texttt{meta\_rprn} & 1.0M & interp & 50.4 & 54.9 & \textbf{build}$^\dagger$ \\
SOSD \texttt{osm\_cellids} & 2M & btree & $-$21.0 & 30.2 & reach., uncapt. \\
SOSD \texttt{fb} & 2M & btree & $-$140.0 & $-$66.7 & reach., uncapt. \\
IDs: \texttt{wiki\_2019t} & 0.80M & btree & $-$267.3 & $-$111.5 & reach., uncapt. \\
\midrule
synth.\ uniform & 2M & interp & 47.2 & 50.2 & unit-dep. \\
synth.\ lognormal & 2M & btree & $-$5.5 & 43.9 & reach., uncapt. \\
synth.\ clustered & 2M & btree & $-$73.6 & $-$11.4 & reach., uncapt. \\
\bottomrule
\end{tabular}
\caption{Learned indexes: percentage of the lookup-cost corridor (tuned non-learned baseline to perfect index) captured by a deployable RMI, under two cost units; both the baseline family and the RMI leaf count are swept. \emph{build} clears the pre-registered 50\% bar on both units. Because reachability holds by construction in this domain (position is $N\cdot\mathrm{CDF}$), the negative rows are not \emph{no corridor} but a fourth category the caching settings do not exhibit, \emph{reachable but uncaptured by this model class}: the corridor is there and a piecewise-linear RMI does not take it. $^\dagger$Clears by 0.4 points, reported as marginal. Against binary search alone the real-key rows would read 91.9, 90.9, 69.7, 40.0 and 22.7\%: the baseline correction is worth up to 337 points of apparent capture.}
\label{tab:index}
\end{table}

\textbf{A substrate-competence control.} A procedure characterised only where it says build has an unmeasured false-negative rate, and the false negative is the expensive error: advising against a component that would have worked is a mistake nobody observes. Two trained cache-management policies were therefore run, both frozen after training on each trace's first half and strictly causal at decision time: a gradient-boosted reuse-distance evictor in the style of LRB \citep{lrb}, and a gradient-boosted admission classifier in the style of Baleen \citep{baleen} with a tuned bypass threshold. Because neither spends speculative bandwidth, neither engages endogenous slack; the trained prefetch-side test is \rung{r6} itself, whose learned timing model does engage it and lands $-40.06$/$-26.01$. What these two arms test instead is whether a learned component of a \emph{different} capability incidentally captures what the procedure declared unreachable, and whether the substrate admits learned wins at all. Each is scored against its own quantity, the \textbf{eviction-management corridor} (Belady over S3-FIFO with no prefetch anywhere), which is distinct from the prefetch corridor of Table~\ref{tab:learnable} and takes different values: $+11.55$, $+12.50$, and $+12.75$ on \texttt{wiki\_2019t}, \texttt{cluster50}, and \texttt{cluster53}. Against those corridors the trained models reach essentially none: the evictor lands $1.83$ points below the baseline on Wikipedia, and the admission classifier captures $-0.0\%$ and $0.0\%$ on the two key-value traces. The latter two are a tuning outcome rather than a capture measurement, and are reported as such: the bypass threshold is tuned on the training half, and tuning selected the degenerate setting, admit everything, so the policy collapses onto S3-FIFO exactly. The failures are informative because the same evictor is competent where the substrate allows, beating S3-FIFO by $+3.45$, $+3.13$, and $+0.99$ points on three traces whose corridors fall below the bar, every interval excluding zero; the failures are therefore properties of the corridors, not of the learners.

\textbf{The recipe.} The instantiations above follow one procedure, stated so it can be run on a new domain.

\vspace{2pt}\noindent\fbox{\begin{minipage}{0.955\columnwidth}\small
\textbf{1.} Fix the deployable baseline by sweeping the strongest non-learned family, tuned against the corridor's interest.
\textbf{2.} Meter every oracle arm to that baseline's measured resource budget, and require Pareto dominance.
\textbf{3.} Restrict the oracle to information a deployable model could hold, and price the restriction inside the replay, not by subtraction.
\textbf{4.} Pre-register thresholds and invariants before the first real run; report every arm the escalation kills.
\textbf{5.} If a corridor survives, price each capability a learner could add with a separate oracle; certify build only if an all-causal combination clears the bar.
\end{minipage}}\vspace{3pt}

\noindent Step 3 is what headroom studies omit, and step 5 is what distinguishes certification from measurement: a corridor only a clairvoyant arm reaches is a trap, and the recipe says so before a model is trained.

\section{Scope and Threats to Validity}
``Learnable'' is defined against frozen, history-based prediction. An online-updating or content-aware predictor could reach cold objects and would raise the surviving corridors, so the measurement verdicts are conservative in that direction; the capture verdicts are not symmetric, since \rung{r4} already grants the widest emission the family supports and \rung{r6}/\rung{r7} fail for timing and redundancy reasons a larger vocabulary does not touch. The two remaining limits on the capture negative, its two-trace scope and the spanning assumption behind Table~\ref{tab:decomp}, are stated where they arise rather than deferred here; the measurement verdicts rest on all twelve traces. The primary operating point is a 1\% cache over 2M-request prefixes, and cache-size and byte-weighted analyses point the same way. Twitter traces are 1-in-10 samples, LLM traces cover one hour of one provider, and intervals cover within-trace variance only. Every error in the baseline inflates corridors, so the negatives are conservative under further scrutiny.

\section{Conclusion}
Measured against a family-tuned baseline at matched bandwidth, most celebrated prefetch headroom in caching is phantom: 14--87\% of gross corridors evaporate across twelve production traces. What survives is real but, priced by pre-registered oracle decompositions, reachable only with foresight, for a reason with a name, endogenous slack, and a general form, rank-learnable is not placeable.

The value of the procedure is that it discriminates: across four settings it returns no corridor (block storage, LLM KV serving), a trap (CDN and key-value prefetching), and a build (learned indexes, on the distributions a piecewise-linear model fits), each decided by replay before training, with a captured corridor and a trained-model competence control bounding the procedure on both sides. Reachability proves necessary but not sufficient, since within the index domain the green light is withheld wherever a tuned classical structure already sits on the reachable ceiling.

The negatives also point somewhere. Endogenous slack penalises volume, not learning, so the productive direction for caching is a predictor that widens reachable coverage without widening low-precision fetch volume, which under the definition used here means content-aware or online-updating prediction, the one class the instrument deliberately excludes and the one whose corridor it would raise. Identifying whether a domain faces a prediction problem or a scheduling problem is exactly what the procedure is for. The procedure and every replay arm used above are provided as anonymized supplementary material, so the next celebrated headroom number can be checked in an afternoon.

\bibliography{p4}

\end{document}
